\documentclass[10pt,a4paper]{amsart}
\usepackage[margin=19mm]{geometry}
\usepackage{amsmath,amssymb,mathtools}
\usepackage[scr=rsfs,cal=euler]{mathalfa}
\usepackage{xfrac}
\usepackage[unicode,hidelinks,bookmarks=false]{hyperref}
\newcommand{\ie}{i.e.\ }
\newcommand{\cf}{cf.\ }
\newcommand{\resp}{respectively\ }
\newcommand{\Subsection}{\subsection}
\providecommand{\nobreakdash}{\nobreak\hbox{-}}
\setlength{\emergencystretch}{2em}
\multlinegap0pt
\usepackage{bm}
\usepackage{bbm}
\usepackage{dsfont}
\usepackage{xspace}
\usepackage{cancel}
\usepackage{mathtools}
\usepackage{amsthm}
\makeatletter
\@ifundefined{lemma}{\newtheorem{lemma}{Lemma}}{}
\@ifundefined{proposition}{\newtheorem{proposition}[lemma]{Proposition}}{}
\makeatother
\newcommand{\cU}{\mathcal{U}}
\newcommand{\cV}{\mathcal{V}}
\newcommand{\cW}{\mathcal{W}}
\newcommand{\supp}{\mathrm{supp}}
\newcommand{\vietm}[1]{\mathrm{V}^\mathrm{m}(#1)}
\title[Vietoris thickenings and complexes of manifolds are homotopy]{Vietoris thickenings and complexes of manifolds are homotopy equivalent}
\author{Henry Adams, Alexandre Karassev, Ziga Virk}
\date{Source: arXiv:2512.23108v1 (CC BY 4.0)}
\begin{document}
\maketitle

\noindent\emph{Source and licence.} Vietoris thickenings and complexes of manifolds are homotopy equivalent. Version arXiv:2512.23108v1 (CC BY 4.0), distributed under the Creative Commons Attribution 4.0 International licence, as recorded in the arXiv licence metadata for this exact version and shown on the abstract page https://arxiv.org/abs/2512.23108v1. Full rights notice: licenses/arxiv-2512.23108-CC-BY-4.0.txt.\par\smallskip

\noindent\textit{Editorial scope.} Counted unit: the Proposition whose source label is \texttt{prop:lsc} in the article (source0.tex lines 844--846, source label \texttt{prop:lsc}), delivered together with the article's own complete original proof of it (source0.tex lines 848--880, one \texttt{proof} environment of 33 lines). No mathematics is written, repaired or re-derived: statement and proof are the article's own text line for line. Required context retained uncounted: lemma:basis3 (lines 730--733); eq:F (lines 818--821). Every definition, display and label of those lines is retained unchanged. Omitted from this package and not redistributed: 1--729, 734--817, 822--843, 847--847, 881--952. Nothing inside a retained excerpt is omitted or altered: every retained body is delivered exactly as the article's lines are written, so no omission receipt of any kind accompanies this package. Citations: the retained excerpts carry no citation command at all, so no bibliography entry of the article is transcribed and no reference is redistributed. Reference adaptation: none was needed and none was made; every reference inside the delivered unit resolves inside it, so no unresolved reference or citation is produced. Printed slips kept literal: no printed word, symbol, hypothesis or number of the counted statement or of its proof was altered. Layout and class: the article's own class is \texttt{amsart} with packages including inputenc, fontenc, amsmath, amsthm, amsfonts, amssymb, url, mathtools, hyperref, enumerate, paralist, tikz, tikz-cd, geometry; the wrapper is the lane's pinned \texttt{amsart} editorial wrapper at 10pt on A4, which restates the theorem-like environments the retained text needs and adds the article's own macro definitions that the retained text needs (\texttt{\textbackslash cU}, \texttt{\textbackslash cV}, \texttt{\textbackslash cW}, \texttt{\textbackslash supp}, \texttt{\textbackslash vietm}), with extra wrapper packages (bm, bbm, dsfont, xspace, cancel, mathtools). The delivered numbering is the unit's own. Assets: no retained span contains a figure environment, a graphics-inclusion command, a PDF-page inclusion, a table or a TikZ picture; no image file is copied into this package, the package declares no asset dependency and no image file is redistributed with it. Licence: version arXiv:2512.23108v1 of this article is distributed under the Creative Commons Attribution 4.0 International licence, as recorded in the arXiv licence metadata for this exact version and shown on the abstract page https://arxiv.org/abs/2512.23108v1. A full rights notice is delivered at licenses/arxiv-2512.23108-CC-BY-4.0.txt. Characters: every retained excerpt is pure ASCII, so the delivered engine, XeLaTeX with its default Latin Modern text font, reports no missing character.
\begin{lemma}
\label{lemma:basis3}
The sets of the form $N(\eta,\cU,\varepsilon)$ form a basis for $\vietm{\cW}$ at $\eta$.
\end{lemma}

\begin{equation}
\label{eq:F}
F(\zeta)=\{\nu\in P(\mu,\cU,\varepsilon)~|~\supp(\nu)\subseteq\supp(\zeta)\}.
\end{equation}

\begin{proposition} \label{prop:lsc}
The multivalued map $F\colon P_{\cU}\cap B(\mu) \to P(X)$ defined in~\eqref{eq:F} is lower semi-continuous.
\end{proposition}

\begin{proof}
To check that $F$ is lower semi-continuous, we must show that for any open subset $S$ of $P(X)$, the set $\{\xi\in P_{\cU}\cap B(\mu)\mid F(\xi)\cap S\ne \varnothing\}$ is open in $P_{\cU}\cap B(\mu)$.
We proceed along the following scheme: for any $\zeta \in P_{\cU}\cap B(\mu)$, any $\nu \in F(\zeta)$, and any arbitrarily small neighborhood $N$ of $\nu$, we define a small neighborhood $N'$ of $\zeta$ such that for each $\xi \in N'$, there exists $\lambda \in F(\xi) \cap N$.

Let
\[\zeta = \sum_{i=1}^n  \sum_{j=1}^{n_i} w^i_j \delta_{y^i_j}+\sum_{t=1}^k w_t \delta_{y_t},\]
where $y^i_j\in U_i$ for all $i,j$, and $y_t\notin \cup\cU$ for all $t=1,2,\dots,k$.
Pick $\nu\in F(\zeta)$.
Next, consider an open neighborhood of $\nu$ which, by Lemma~\ref{lemma:basis3}, may be assumed to be of the form $N=N(\nu,\widetilde \cU,\alpha)$.
Moreover, if 
$\cV=\{V^i_j\mid i=1,\dots,n,\, j=1,\dots, n_i\}$ is a collection of 
%$\{V^i_j\mid i=1,2,\dots ,n, j=1,2,\dots, n_i\}$ are 
disjoint open neighborhoods of the points $y^i_j$ such that $V^i_j\subseteq U_i$ for all $i$, 
and if $\cV'=\{V^i_j\mid y^i_j \in \supp \nu\}$, then
we may further assume that 
$N=N(\nu,\cV',\alpha)$.
%$N$ has the form $N(\nu, \{V^i_j\mid y^i_j \in \supp \nu\},\alpha)$.
Choose $\beta >0$ so that $\beta < w^i_j$ for all $i,j$, and so that
$N' \coloneqq N(\zeta,\cV,\beta)$
%\[N' \coloneqq N(\zeta, \{V^i_j\mid \ i=1,2,\dots,n, j=1,2,\dots ,n_i\},\beta)\]
is contained in $P_{\cU}\cap B(\mu)$, the domain of $F$.
Consider $\xi \in N'$.
Note that the choice of $\beta$ implies that $\supp \xi \cap V^i_j \ne \varnothing$ for all $i,j$.
For each $i=1,2,\dots ,n$ and $j=1,2,\dots, n_i$, pick $z^i_j\in \supp \xi \cap V^i_j$, and let 
\[\lambda = \sum_{
V^i_j\in \cV'
%\nu(V^i_j)\ne 0
} \nu (V^i_j) \delta_{z^i_j}.\]
It is easy to see that $\lambda (U_i) = \nu(U_i)$ for $i=1,2,\dots ,n$, and that $\supp\lambda \subseteq \supp \xi$.
Therefore $\lambda \in F(\xi)$.
Moreover, $\lambda (V^i_j) = \nu (V^i_j)$ for all $V^i_j\in \cV'$, and hence $\lambda \in N$.
Thus, $\lambda \in F(\xi)\cap N$, as required.
\end{proof}

\end{document}
